\section{Independent Evaluation Methodology}
\label{sec:independent_evaluation}

The original and augmented POD bases are compared using a
separate test dataset that is independent of both training
datasets. This evaluation is designed to assess whether the
additional snapshots improve the predictive accuracy of the
reduced-order model at parameter combinations that were not
used during the construction of either basis.

Both bases are evaluated using the same full-order reference
solutions, parameter combinations, and retained mode counts.
This paired evaluation makes it possible to attribute
differences in ROM accuracy to the choice of POD basis rather
than to differences in the evaluation cases.

\subsection{Generation of the final test set}

A second set of 100 parameter combinations is generated using
Latin hypercube sampling (LHS) over the prescribed parameter
domain:
\[
1\leq\mu\leq100,
\qquad
0.01\leq Br\leq2.
\]

Unlike the first LHS dataset, which was used to screen the
original basis and identify high-error regions, this dataset
is reserved for the final comparison between the original
and augmented bases. It is generated only after the
additional training snapshots have been selected and the
augmented POD basis has been constructed.

To maintain the independence of the test set, the parameter
combinations used in both training datasets are first
collected. These include the original 100 training cases
and the 15 additional cases introduced through local
refinement. The corresponding parameter combinations are
stored as training keys and used to filter the LHS samples.

The test-grid generation procedure follows these steps:
\begin{enumerate}
    \item Generate candidate parameter combinations using
    a Latin hypercube sampler over the prescribed domain.
    \item Load the parameter keys associated with the
    original and augmented training datasets.
    \item Reject candidate combinations that duplicate
    parameter combinations already used in either training
    dataset.
    \item Retain unique parameter combinations to construct
    the final test grid.
    \item Save the accepted combinations in a CSV file.
\end{enumerate}

After obtaining the 100 test parameter combinations, the
full-order solver is executed for each case. The resulting
converged temperature fields are saved and associated with
their parameter values in a test-case summary file.

The final evaluation dataset therefore contains 100
parameter combinations and their corresponding full-order
reference temperature fields. None of these parameter
combinations is included in the original or augmented
POD training datasets.

\subsection{Paired comparison of the two POD bases}

The original and augmented POD bases are evaluated using
exactly the same final test cases. For each test parameter
combination, the corresponding full-order temperature
solution is used as the common reference for both ROM
evaluations.

The comparison considers the retained mode counts
\[
\boxed{
r\in\{1,2,3,4,5,6,8,10\}.
}
\]

For every mode count, the following procedure is performed:
\begin{enumerate}
    \item Load the original and augmented POD bases.
    \item Select the first $r$ modes from each basis.
    \item Evaluate the original-basis ROM for each test case.
    \item Evaluate the augmented-basis ROM for the same
    test case and retained mode count.
    \item Reconstruct both temperature fields and compare
    them against the same full-order reference solution.
    \item Store the case-wise errors and pair the original
    and augmented results using their test-case identifiers.
\end{enumerate}

Since both bases are evaluated at every mode count for all
100 test cases, the comparison involves
\[
100\times8\times2=1600
\]
ROM evaluations: 800 evaluations using the original basis
and 800 using the augmented basis.

The paired evaluation is important because it keeps the
parameter combinations and reference solutions unchanged
while changing only the POD basis. Consequently, the
differences in error can be examined on a case-by-case
basis rather than inferred solely from two independently
calculated error distributions.

\subsection{Comparison statistics}

The relative $L_2$ error defined in
Eq.~\eqref{eq:relative_l2_error} is used as the primary
measure for comparing the two bases. For each basis and
retained mode count, the mean, median, 95th percentile,
and maximum relative $L_2$ errors are calculated over
the 100 test cases.

The mean error describes the overall average performance,
while the median represents the central value of the error
distribution. The 95th percentile provides information
about the upper tail of the distribution without depending
on only the single worst case. The maximum error identifies
the most inaccurate test case.

These statistics are reported separately for the original
and augmented bases in Table~\ref{tab:independent_error_statistics}.

\begin{table}[htbp]
    \centering
    \caption{Relative $L_2$ error statistics for the original
    and augmented POD bases over the 100 independent test cases.}
    \label{tab:independent_error_statistics}
    \scriptsize
    \begin{tabular}{ccccccc}
        \toprule
        \textbf{Basis} &
        \textbf{Modes} &
        \textbf{Cases} &
        \textbf{Mean} &
        \textbf{Median} &
        \textbf{95th percentile} &
        \textbf{Maximum} \\
        \midrule
        Original  & 1  & 100 & 0.401768 & 0.437246 & 0.620615 & 0.634832 \\
        Original  & 2  & 100 & 0.176187 & 0.166973 & 0.371498 & 0.411629 \\
        Original  & 3  & 100 & 0.028506 & 0.024564 & 0.062156 & 0.081035 \\
        Original  & 4  & 100 & 0.008127 & 0.007205 & 0.026946 & 0.044410 \\
        Original  & 5  & 100 & 0.002725 & 0.001065 & 0.018959 & 0.030992 \\
        Original  & 6  & 100 & 0.001571 & 0.000189 & 0.013137 & 0.026048 \\
        Original  & 8  & 100 & 0.000876 & $3.00\times10^{-6}$ & 0.006426 & 0.018205 \\
        Original  & 10 & 100 & 0.000577 & $3.31\times10^{-8}$ & 0.003450 & 0.013008 \\
        \midrule
        Augmented & 1  & 100 & 0.384505 & 0.423167 & 0.594550 & 0.608058 \\
        Augmented & 2  & 100 & 0.193602 & 0.163706 & 0.435441 & 0.468897 \\
        Augmented & 3  & 100 & 0.036678 & 0.035301 & 0.079193 & 0.095283 \\
        Augmented & 4  & 100 & 0.013629 & 0.013465 & 0.028654 & 0.037714 \\
        Augmented & 5  & 100 & 0.003841 & 0.004188 & 0.007181 & 0.009805 \\
        Augmented & 6  & 100 & 0.000930 & 0.000758 & 0.001740 & 0.005040 \\
        Augmented & 8  & 100 & 0.000141 & 0.000123 & 0.000241 & 0.001747 \\
        Augmented & 10 & 100 & $1.74\times10^{-5}$ & $1.30\times10^{-5}$ & $3.07\times10^{-5}$ & 0.000401 \\
        \bottomrule
    \end{tabular}
\end{table}

In addition to the aggregate statistics, the comparison
examines the difference in relative $L_2$ error for each
paired test case. Define the case-wise error change as
\begin{equation}
    \boxed{
    \Delta\epsilon_k
    =
    \epsilon_{\mathrm{rel},k}^{\mathrm{aug}}
    -
    \epsilon_{\mathrm{rel},k}^{\mathrm{orig}}.
    }
    \label{eq:casewise_error_change}
\end{equation}
A negative value indicates that the augmented basis
produces a lower error for that test case, whereas a
positive value indicates that its error is higher.

The mean and median of $\Delta\epsilon_k$, together
with its 5th and 95th percentiles, are calculated for
each retained mode count. The percentage of test cases
for which the augmented basis produces a lower error
is also evaluated:
\begin{equation}
    \boxed{
    P_{\mathrm{improved}}
    =
    \frac{
    \displaystyle\sum_{k=1}^{N_t}
    \mathbb{I}
    \left(
    \epsilon_{\mathrm{rel},k}^{\mathrm{aug}}
    <
    \epsilon_{\mathrm{rel},k}^{\mathrm{orig}}
    \right)
    }{N_t}
    \times100,
    }
    \label{eq:percentage_improved}
\end{equation}
where $N_t=100$ and $\mathbb{I}$ is the indicator
function.

The percentage reduction in the mean relative $L_2$
error is defined as
\begin{equation}
    \boxed{
    R_{\mathrm{mean}}
    =
    \frac{
    \bar{\epsilon}_{\mathrm{orig}}
    -
    \bar{\epsilon}_{\mathrm{aug}}
    }{
    \bar{\epsilon}_{\mathrm{orig}}
    }\times100.
    }
    \label{eq:mean_error_reduction}
\end{equation}
An analogous expression is used for the maximum relative
$L_2$ error:
\begin{equation}
    R_{\max}
    =
    \frac{
    \epsilon_{\max,\mathrm{orig}}
    -
    \epsilon_{\max,\mathrm{aug}}
    }{
    \epsilon_{\max,\mathrm{orig}}
    }\times100.
    \label{eq:maximum_error_reduction}
\end{equation}
Positive values indicate a reduction in error with the
augmented basis, whereas negative values indicate an
increase.

Reporting only the mean error would be insufficient for
this comparison. An augmented basis may substantially
reduce the mean error through improvements in a relatively
small number of cases, while producing larger errors in
many other cases. The case-wise difference and percentage
of improved cases provide additional information about
whether the enrichment benefits the test set consistently.

\subsection{Results of the paired basis comparison}

The paired comparison results for the eight retained
mode counts are presented in
Table~\ref{tab:paired_basis_comparison}. The table reports
the mean and maximum relative $L_2$ errors for each basis,
their percentage reductions, and the percentage of
individual test cases improved by the augmented basis.

\begin{table}[htbp]
    \centering
    \caption{Paired comparison of original and augmented
    POD bases over 100 independent test cases.}
    \label{tab:paired_basis_comparison}
    \scriptsize
    \setlength{\tabcolsep}{3pt}
    \resizebox{\textwidth}{!}{%
        \begin{tabular}{cccccccc}
            \toprule
            \textbf{Modes} &
            \textbf{Mean, orig.} &
            \textbf{Mean, aug.} &
            \textbf{Mean reduction (\%)} &
            \textbf{Max., orig.} &
            \textbf{Max., aug.} &
            \textbf{Max. reduction (\%)} &
            \textbf{Improved cases (\%)} \\
            \midrule
            1  & 0.401768 & 0.384505 & 4.2967   & 0.634832 & 0.608058 & 4.2175   & 100 \\
            2  & 0.176187 & 0.193602 & -9.8843  & 0.411629 & 0.468897 & -13.9125 & 33 \\
            3  & 0.028506 & 0.036678 & -28.6650 & 0.081035 & 0.095283 & -17.5825 & 32 \\
            4  & 0.008127 & 0.013629 & -67.6905 & 0.044410 & 0.037714 & 15.0773  & 28 \\
            5  & 0.002725 & 0.003841 & -40.9432 & 0.030992 & 0.009805 & 68.3639  & 22 \\
            6  & 0.001571 & 0.000930 & 40.7665  & 0.026048 & 0.005040 & 80.6523  & 22 \\
            8  & 0.000876 & 0.000141 & 83.9121  & 0.018205 & 0.001747 & 90.4018  & 11 \\
            10 & 0.000577 & $1.74\times10^{-5}$ & 96.9841 &
            0.013008 & 0.000401 & 96.9175 & 20 \\
            \bottomrule
        \end{tabular}%
    }
\end{table}

The case-wise error changes are summarized in
Table~\ref{tab:casewise_error_changes}. Negative values
represent improvements produced by the augmented basis.

\begin{table}[htbp]
    \centering
    \caption{Statistics of the paired case-wise relative
    $L_2$ error changes,
    $\Delta\epsilon=\epsilon_{\mathrm{aug}}-\epsilon_{\mathrm{orig}}$.}
    \label{tab:casewise_error_changes}
    \small
    \begin{tabular}{cccccc}
        \toprule
        \textbf{Modes} &
        \textbf{Mean} &
        \textbf{Median} &
        \textbf{5th percentile} &
        \textbf{95th percentile} &
        \textbf{Improved cases (\%)} \\
        \midrule
        1  & -0.017260 & -0.017340 & -0.026400 & -0.004820 & 100 \\
        2  &  0.017415 &  0.041945 & -0.080070 &  0.077803 & 33 \\
        3  &  0.008171 &  0.011487 & -0.027850 &  0.031049 & 32 \\
        4  &  0.005501 &  0.007381 & -0.014240 &  0.021053 & 28 \\
        5  &  0.001116 &  0.002426 & -0.013540 &  0.005900 & 22 \\
        6  & -0.000640 &  0.000445 & -0.011450 &  0.001546 & 22 \\
        8  & -0.000730 &  0.000082 & -0.006310 &  0.000193 & 11 \\
        10 & -0.000560 &  0.00000833 & -0.003440 &  0.0000271 & 20 \\
        \bottomrule
    \end{tabular}
\end{table}

For one retained mode, the augmented basis improves all
100 test cases. The mean relative $L_2$ error decreases
by 4.30\%, and the maximum error decreases by 4.22\%.

The results are different for two and three modes. At
both ranks, the augmented basis produces larger mean
and maximum errors. Only 33\% and 32\% of the test
cases, respectively, have lower errors with the
augmented basis.

For four and five modes, the maximum errors decrease
by 15.08\% and 68.36\%, respectively. However, the
mean errors increase by 67.69\% and 40.94\%, and
only 28\% and 22\% of the test cases improve. Thus,
the reduction in the worst-case error does not imply
an overall improvement across the test set.

At six modes, the augmented basis reduces the mean
relative $L_2$ error by 40.77\% and the maximum error
by 80.65\%. Nevertheless, only 22\% of individual
test cases improve. The mean and median case-wise
error changes also have opposite signs. This indicates
that substantial improvements in a relatively small
number of cases compensate for error increases in
a larger number of cases.

The difference is even more pronounced at eight
and ten modes. At eight modes, the mean error
decreases by 83.91\% and the maximum error decreases
by 90.40\%, but only 11\% of the individual cases
improve. At ten modes, the mean and maximum errors
decrease by approximately 96.98\% and 96.92\%,
respectively, while only 20\% of the cases improve.

These results show why the mean error alone cannot
describe the effect of snapshot enrichment. At higher
basis dimensions, the augmented basis produces
substantial improvements in the mean and maximum
errors, but the case-wise statistics show that these
improvements are not distributed uniformly across
the test set. For several mode counts, the median
case-wise error change is positive even when the
mean change is negative.

Additional error measures are reported in
Table~\ref{tab:additional_error_metrics}. These include
the mean and median RMSE, the mean of the case-wise
maximum absolute errors, and the largest absolute
error across all test cases.

\begin{table}[htbp]
    \centering
    \caption{Additional error measures and mean ROM
    computation time for the two POD bases.}
    \label{tab:additional_error_metrics}
    \scriptsize
    \setlength{\tabcolsep}{3pt}
    \resizebox{\textwidth}{!}{%
        \begin{tabular}{ccccccc}
            \toprule
            \textbf{Basis} &
            \textbf{Modes} &
            \textbf{Mean RMSE} &
            \textbf{Median RMSE} &
            \textbf{Mean max. abs. error} &
            \textbf{Maximum abs. error} &
            \textbf{Mean ROM time (s)} \\
            \midrule
            Original  & 1  & 0.815932 & 0.800746 & 1.535463 & 5.172093 & 0.000307 \\
            Original  & 2  & 0.378037 & 0.302243 & 0.795553 & 3.403278 & 0.000402 \\
            Original  & 3  & 0.074669 & 0.040121 & 0.171397 & 1.960261 & 0.000366 \\
            Original  & 4  & 0.026578 & 0.011539 & 0.068964 & 1.086724 & 0.000386 \\
            Original  & 5  & 0.012879 & 0.001724 & 0.036530 & 0.750323 & 0.000408 \\
            Original  & 6  & 0.009093 & 0.000303 & 0.027099 & 0.636573 & 0.000415 \\
            Original  & 8  & 0.005597 & $4.91\times10^{-6}$ & 0.018300 & 0.516899 & 0.000461 \\
            Original  & 10 & 0.003786 & $4.71\times10^{-8}$ & 0.013759 & 0.434450 & 0.000523 \\
            \midrule
            Augmented & 1  & 0.774306 & 0.741124 & 1.498850 & 4.911156 & 0.000331 \\
            Augmented & 2  & 0.380486 & 0.286150 & 0.748525 & 2.356156 & 0.000417 \\
            Augmented & 3  & 0.080750 & 0.061802 & 0.175861 & 1.057812 & 0.000418 \\
            Augmented & 4  & 0.029537 & 0.022399 & 0.070843 & 0.541289 & 0.000431 \\
            Augmented & 5  & 0.009004 & 0.006955 & 0.023487 & 0.296553 & 0.000461 \\
            Augmented & 6  & 0.002595 & 0.001311 & 0.008216 & 0.193327 & 0.000473 \\
            Augmented & 8  & 0.000475 & 0.000177 & 0.001579 & 0.073051 & 0.000544 \\
            Augmented & 10 & $7.24\times10^{-5}$ & $2.26\times10^{-5}$ & 0.000221 & 0.014533 & 0.000637 \\
            \bottomrule
        \end{tabular}%
    }
\end{table}

The absolute-error measures show the same general trend
at higher basis dimensions: the augmented basis reduces
the mean RMSE and the maximum absolute error for six,
eight, and ten modes. Its mean ROM computation time is
slightly higher for every mode count. This is consistent
with the additional projection and reduced-system
operations associated with the larger augmented
snapshot basis. These times describe the ROM evaluation
itself and do not include the cost of generating the
full-order training and test datasets.

Overall, the augmented basis does not provide a uniform
improvement across all retained mode counts and test
cases. Its performance is worse on average for two,
three, four, and five modes, although the maximum
error decreases at four and five modes. At six, eight,
and ten modes, it produces substantially lower mean
and maximum errors, but the percentage of individual
cases improved remains relatively small.

The results therefore indicate that the effect of
snapshot enrichment depends on the retained basis
dimension and the parameter combination being
evaluated. The improvement in aggregate error
statistics should be interpreted together with the
paired case-wise differences.

\subsection{Computational procedure and code implementation}

The final independent evaluation is implemented through
two scripts in the \texttt{F05\_refinement} directory.

\begin{itemize}

    \item \texttt{f06\_generate\_independent\_test\_set.py}:
    Generates the final independent LHS test dataset.
    Its main function,
    \texttt{run\_independent\_test\_set\_generation},
    loads the parameter keys from both the original
    and augmented training datasets. The function
    \texttt{load\_training\_parameter\_keys} combines
    these keys to exclude previously used parameter
    combinations.

    The function \texttt{generate\_test\_parameter\_grid}
    generates candidate samples using
    \texttt{scipy.stats.qmc.LatinHypercube}, filters
    duplicate and previously used parameter
    combinations, and returns the accepted test grid.
    The resulting parameter combinations are saved
    in a CSV file. The main function then processes
    each accepted case using \texttt{run\_fom\_case},
    saves the full-order temperature field, and
    records the associated simulation information.

    \item \texttt{f07\_compare\_original\_and\_augmented\_bases.py}:
    Evaluates and compares the original and augmented
    POD bases using the independent test dataset.
    The main function, \texttt{run\_basis\_comparison},
    loads the full-order test cases and both POD bases,
    evaluates the ROMs over the selected mode counts,
    constructs paired case-wise comparisons, and
    generates the corresponding error visualizations.

    The function \texttt{evaluate\_rom} performs an
    individual ROM evaluation. It constructs the
    full-order operator components, projects them
    onto the selected POD basis, assembles the
    reduced system, solves for the modal coefficients,
    and reconstructs the temperature field.

    The function \texttt{run\_rom\_evaluations}
    applies this procedure to every test case, both
    POD bases, and all selected mode counts. Each
    result is saved with its corresponding test-case
    identifier and basis information.

    The function \texttt{build\_comparison\_summary}
    calculates the aggregate error statistics for
    each basis and mode count. The function
    \texttt{build\_paired\_comparison} associates
    the original and augmented results for each
    identical test case. The function
    \texttt{build\_improvement\_summary} then
    calculates the case-wise error changes and
    improvement statistics.

    Finally, the plotting routines
    \texttt{plot\_error\_comparison},
    \texttt{plot\_error\_reduction}, and
    \texttt{plot\_error\_maps} generate visual
    comparisons of the error distributions, the
    changes introduced by augmentation, and their
    variation across the parameter domain.

\end{itemize}

The evaluation workflow can be summarized as
\[
\boxed{
\begin{gathered}
\text{Independent LHS parameter set}\\
\downarrow\\
\text{Full-order reference solutions}\\
\downarrow\\
\begin{array}{cc}
\text{Original POD basis} &
\text{Augmented POD basis}\\
\downarrow & \downarrow\\
\text{Original ROM} & \text{Augmented ROM}
\end{array}\\
\downarrow\\
\text{Paired error comparison}\\
\downarrow\\
\text{Aggregate and case-wise statistics}
\end{gathered}
}
\]

Since both bases use the same independent test cases and
full-order reference fields, the resulting comparison
provides a consistent assessment of the effect of
snapshot enrichment on ROM accuracy.
